\section{Exercises}\label{sec:04-propositions-and-proofs:08-exercises:1}

\textbf{Key points.} \texttt{Prop} is the type of statements, and a proof is just an ordinary term of that type (Curry--Howard). \texttt{∧}/\texttt{∨}/\texttt{→}/\texttt{∀}/\texttt{∃} each have their own introduction and elimination rule, mirrored directly by Lean syntax, the anonymous constructor, \texttt{Or.inl}/\texttt{Or.inr}/\texttt{match}, \texttt{fun}/ application, and \texttt{fun}/anonymous-constructor again. \texttt{rfl} proves equality by checking both sides reduce to the same normal form.

\begin{enumerate}
\def\labelenumi{\arabic{enumi}.}
\tightlist
\item
  Prove \texttt{theorem and\_\allowbreak{}comm\_\allowbreak{}ex \{P Q : Prop\} (h : P ∧ Q) : Q ∧ P}.
\item
  Prove \texttt{theorem or\_\allowbreak{}comm\_\allowbreak{}ex \{P Q : Prop\} (h : P ∨ Q) : Q ∨ P} (hint, use \texttt{Or.elim} or pattern matching with \texttt{match}).
\item
  Prove \texttt{theorem exists\_\allowbreak{}gt\_\allowbreak{}zero : ∃ n : Nat, n > 0}.
\item
  \texttt{P ∧ Q} has a single constructor, the anonymous pair; \texttt{P ∨ Q} has two, \texttt{Or.inl} and \texttt{Or.inr}. State precisely what data a proof of each connective must carry, and show that this difference in data forces the difference in constructor count.
\item
  Show that \texttt{rfl} proves \texttt{theorem nat\_\allowbreak{}add\_\allowbreak{}zero (n : Nat) : n + 0 = n} but cannot prove \texttt{∀ n, 0 + n = n} by the same route. Identify the clause of \texttt{Nat.add}'s recursion responsible for the asymmetry.
\item
  \texttt{exists\_\allowbreak{}prime\_\allowbreak{}gt\_\allowbreak{}three} (Section 6) was proved with one witness and a decided fact about it. Explain why no finite number of such witness-and-proof pairs could ever establish \texttt{∀ n, n > 0}, and name the kind of argument that Chapter 5 introduces to close this gap.
\end{enumerate}

Solutions, \hyperref[sec:15-appendix-solutions:04-chapter-4:1]{Appendix, Chapter 4}.

